\chapter{Conclusion and Future Scope}
\label{ch:concl}

\section{Conclusion and Future Perspectives}

\subsection{Project Summary and Key Contributions}
Verif-AI is a working system for face-based deepfake detection and forensic analysis. It locates faces with MTCNN, scores them with a two-stream
network (an RGB EfficientNet-B0 and an EfficientNet-B0 that reads a fixed three-kernel SRM noise residual), explains the score with Grad-CAM heat-maps,
summarises videos with face tracks, timelines and stored evidence frames, and produces a PDF forensic report. It is delivered as a tested FastAPI
backend and a responsive React dashboard in which uploads are deleted immediately and results can be deleted by the user. Beyond the software, the
project contributes an evaluation record that treats failures as results: the chance-level behaviour of the production model on genuine video was
reproduced, traced to an image-quality shortcut and partly repaired; a per-quality-bin acceptance gate was designed and validated with synthetic
controls; and an experiment with modern-generator data was run under a protocol and a promotion rule that were fixed before training.

\subsection{Final Performance Assessment}
The honest summary of the measured performance is the following. The production model v2 reaches accuracy 0.875 and ROC-AUC 0.941 on held-out identities
of DeepFakeFace with generators seen in training, but only 0.746 and 0.849 on an unseen text-to-image generator, and it is at chance level (accuracy
0.483) on held-out Celeb-DF video. The model with Celeb-DF data and quality augmentation (Model 4), the previous default and now an alternative, reaches 0.980 video accuracy on
held-out Celeb-DF identities --- within-dataset evidence for one face-swap method --- but detects fewer DeepFakeFace fakes than v2. The E3 models detect held-out
fakes of the modern generator families they were trained on almost perfectly, but this is not generator-disjoint evidence and the models (E3A, and the source-balanced E4) did not meet the
pre-registered promotion rule. No result supports a claim of 95\% accuracy for modern AI-generated images or videos in general, and none supports a
claim about AI-generated video.

\subsection{Critical Reflections and Limitations}
The central lesson is methodological. The prototype's very high reported accuracy could not be reproduced on identity-disjoint data, an
unseen generator exposed a large drop, and a cue as simple as sharpness explained much of the apparent skill on the original dataset.
Scores are not probabilities, Grad-CAM maps are explanations of a model and not proof, and the system depends on face detection and on the data it was
trained on. The public demonstration runs on free hosting with a smaller, older model, without heat-maps and without persistent storage, and has no
authentication or rate limiting; it is a demonstration, not a production service. Section~\ref{sec:limits} lists the limitations in detail.

\subsection{Future Perspectives}
\begin{itemize}
\item \textbf{Independent evaluation data.} Assemble labelled test sets with \emph{generator-disjoint} and \emph{source-disjoint} splits, genuine
phone and webcam photographs and video, and, if licences allow, labelled AI-generated video; only then can claims about modern generators be tested.
\item \textbf{Shortcut-robust training.} Balance training data across quality bins and sources, supervise each stream separately, and report the per-bin
gate for every candidate; the source-balanced E4 experiment (Section~\ref{sec:e4}) is a first step in this direction.
\item \textbf{Calibration and uncertainty.} Calibrate scores on held-out data of the intended use and report calibrated uncertainty.
\item \textbf{Stronger explanation.} Store heat-maps for every analysed frame together with their registered position so that a true overlay is possible for video.
\item \textbf{Production hardening.} Authentication, rate limiting, persistent encrypted storage, monitoring and a hosting plan with enough memory for Grad-CAM and longer videos.
\item \textbf{Broader inputs.} Audio-visual and temporal models for video, and analysis beyond faces, are outside the scope of this work.
\end{itemize}

\subsection{Closing Remarks}
Verif-AI shows how a deepfake detector can be built as an inspectable, privacy-conscious tool and, more importantly, how its claims should be tested and bounded.
A detector is only as trustworthy as the evidence about where it fails, and this report has tried to give that evidence the same weight as the architecture.
